\section{源码等价的 A--H 执行模型}
\label{sec:evpt-source-equivalent}

本章按公开入口分别记录实现，禁止把不同入口拼成一个不存在的统一优化模型。所有严格模式均保留源码
拒绝语义：\texttt{simulate\_ev\_power\_traffic} 的路线分配、充电调度和电网验证不是一个全局认证
程序，\texttt{require\_exact\_mathematical\_model=true} 时直接返回不支持。依据：
\srcpath{src/ev\_power\_traffic/ev\_power\_traffic\_simulation.cpp:61--95}。

\subsection{A：分解仿真入口}
入口为 \texttt{simulate\_ev\_power\_traffic}。它先生成或读取候选路径，再按
\texttt{assignment\_model} 执行路径分配，生成充电会话并逐时验证电网。最终
\texttt{feasible} 的源码谓词是：系统最优求解请求必须求解成功；若禁止未服务出行，则未服务车辆
不超过 $10^{-7}$；若禁止未服务充电，则未服务电量不超过 $10^{-6}$；每步发电容量必须合格，且
已经运行的潮流必须收敛。依据：
\srcpath{src/ev\_power\_traffic/ev\_power\_traffic\_simulation.cpp:7--41}。

该入口不是单体数学规划，因此本章不为 A 添加联合目标函数。站点负荷写回只执行
\begin{equation}
 p_s^{model}=p_s^{EV},\qquad q_s^{model}=q_s^{EV},\qquad
 utilization_s=\begin{cases}|p_s^{EV}|/P_s^{max},&P_s^{max}>10^{-9},\\0,&\text{否则。}\end{cases}
\end{equation}
依据：\srcpath{src/ev\_power\_traffic/ev\_power\_traffic\_simulation.cpp:46--58}。

\subsection{B/F：CTM 传播与 DUE--VI}
对路段 $a$，若长度和 CTM 步长有效，元胞数为
\begin{equation}
 M_a=\begin{cases}
 \texttt{n\_cells\_per\_link},&\text{该选项}>0,\\
 \max(1,\lfloor L_a/(v_a^f\Delta t_c)\rfloor),&\text{否则，}
 \end{cases}
 \qquad\delta_a=L_a/M_a.
\end{equation}
无效长度或步长时固定 $M_a=1$ 且 $\delta_a=\max(L_a,10^{-9})$。容量非正时取
$10^4$ veh/h。拥堵波速按
\begin{equation}
 k^{jam}=\begin{cases}N^{jam}/\delta,&N^{jam}>10^{-9},\\2q^{max}/v^f,&\text{否则，}\end{cases}
 \quad k^{crit}=q^{max}/v^f,
 \quad w=\begin{cases}q^{max}/(k^{jam}-k^{crit}),&k^{jam}>k^{crit}+10^{-9},\\
 w^{fallback},&\text{否则。}\end{cases}
\end{equation}
依据：\srcpath{src/ev\_power\_traffic/ctm\_propagation.cpp:40--103}。

元胞发送、接收和内部边界流严格为
\begin{align}
 S_{a,m,k}&=\Delta t_c\min(n_{a,m,k}v_a^f/\delta_a,q_{a,k}^{max}),\\
 R_{a,m,k}&=\Delta t_c\min(q_{a,k}^{max},
 w_a\max(0,N_a^{jam}-n_{a,m,k})/\delta_a),\\
 y_{a,m,k}&=\min(S_{a,m,k},R_{a,m+1,k}),\\
 n_{a,m,k+1}&=\max(0,n_{a,m,k}+y_{a,m-1,k}-y_{a,m,k}).
\end{align}
路径专属元胞开启时，跨边界的各路径流按上游旧占有量比例分配，而非独立求最小值。节点合流/分流
还有两阶段预算缩放，不能用上述内部元胞式替代；权威实现为
\srcpath{src/ev\_power\_traffic/ctm\_propagation.cpp:401--752}。内部元胞依据：
\srcpath{src/ev\_power\_traffic/ctm\_propagation.cpp:110--127}、
\srcpath{src/ev\_power\_traffic/ctm\_propagation.cpp:313--385}。

F 入口为 \texttt{solve\_ctm\_due\_vi}。精确模型开关为真时直接拒绝，因为结果只给固定点残差，
不是全局 VI/MPEC 证明。仿真步与 CTM 步的倍率是
\begin{equation}
 R=\max(1,\operatorname{round}(\Delta t_{sim}/\Delta t_c)),\qquad T_c=T_{sim}R.
\end{equation}
初始流在每个需求的出发步集合和可行候选路径间等分。混合车流开启时，送入 CTM 的流为
\begin{equation}
 x^{combined}_{k,r}=PCE_{EV}x^{EV}_{k,r}+PCE_{ICV}x^{ICV}_{k,r}.
\end{equation}
依据：\srcpath{src/ev\_power\_traffic/ctm\_due\_vi.cpp:66--106}、
\srcpath{src/ev\_power\_traffic/ctm\_due\_vi.cpp:166--236}。

每轮全有或全无方向为 $\mathbf d=\mathbf y-\mathbf x$。若
\texttt{line\_search\_steps}>0，代码在 $[0,1]$ 上用黄金分割最小化 CTM 重放得到的
$TSTT=\sum x_{k,r}t_{k,r}$；否则
\begin{equation}
 \alpha_i=\frac1{i+1},\qquad\mathbf x\leftarrow\mathbf x+\alpha_i\mathbf d.
\end{equation}
这不是旧文档所写的 $1/i$。收敛先按 Wardrop 相对间隙小于
\texttt{convergence\_tol} 判定；最终若 VI 证书可用，报告间隙还取既有报告、EV/ICV 间隙与
归一化 VI 间隙的最大值。混合车流时 VI 证书只覆盖 EV 子博弈。依据：
\srcpath{src/ev\_power\_traffic/ctm\_due\_vi.cpp:239--430}。

\subsection{C：价格协调，不是单体最优}
C 每轮先运行 CTM--DUE，再逐时把站点有功以 $p_s^{EV}/1000$ MW 加入母线负荷并求 DC--OPF。
OPF 收敛时，站价更新为
\begin{equation}
 \pi^{new}_{s,k}=\alpha\frac{LMP_{bus(s),k}}{1000}+(1-\alpha)\pi^{old}_{s,k},
 \qquad\alpha=\operatorname{clamp}(\texttt{price\_update\_step},0,1).
\end{equation}
OPF 失败的步保留旧价。收敛仅指
\begin{equation}
 \max_{s,k}|\pi^{new}_{s,k}-\pi^{old}_{s,k}|<
 \texttt{price\_convergence\_tol};
\end{equation}
它不证明社会福利最优。报告福利按代码后处理为
$W=B_{EV}-C_{gen}-VOT\cdot TSTT$。禁用 DC--OPF 时固定价格且只运行一轮，结果被置为收敛；
此时报告生成成本为零而不是静态估计值。依据：
\srcpath{src/ev\_power\_traffic/ctm\_joint\_welfare.cpp:237--355}。

\subsection{D：按模式分派的联合程序}
入口 \texttt{solve\_joint\_optimizer} 根据 \texttt{JointOptimizerMode} 分派到不同模型：基础
LP/MILP、有限路线 BPR NLP/MPEC、认证动态 MILP 或全联合 LTM--PWL MILP。它们的变量、约束和
证书不同，不能共享一条目标式。基础模型的权威装配为
\srcpath{src/ev\_power\_traffic/joint\_optimizer.cpp:1652--2644}；非线性模式为
\srcpath{src/ev\_power\_traffic/joint\_optimizer\_full\_nlp.cpp:1760--1915}。

结果字段已区分 \texttt{proven\_optimal}、\texttt{local\_optimum\_certificate}、
\texttt{global\_optimum\_certificate}、\texttt{travel\_times\_are\_endogenous}、
\texttt{wardrop\_complementarity\_enforced} 和 \texttt{dcopf\_coupling\_enforced}。
文档不得把局部 KKT、PWL 近似或一次 CTM 预运行的外生阻抗改写为全局联合最优。

\subsection{E 与 E--LTM}
\texttt{solve\_ctm\_so\_lp} 和 \texttt{solve\_ltm\_so\_lp} 是不同的系统最优 LP 入口，
分别在 \srcpath{src/ev\_power\_traffic/ctm\_so\_lp.cpp:15} 与
\srcpath{src/ev\_power\_traffic/ltm\_so\_lp.cpp:39} 装配。前者使用时间展开 CTM 节点/路径连续性；后者
使用累积计数 LTM。由于两份装配的行结构不同，本章不使用 Beckmann 积分式替代其中任一实现。

\subsection{G：约化选址定容 MILP}
严格模型开关为真时 G 直接拒绝，状态说明该 MILP 使用目标/动态近似。实际站点变量为二进制
$z_s$ 和整数 $B_s\in[0,B^{max}]$，并加入
\begin{equation}
 B_s-B^{max}z_s\le0.
\end{equation}
每个有效需求建立
\begin{equation}
 \sum_{r\in R_d}h_{d,r}+u_d=D_d,
\end{equation}
其中路径流的初始目标系数为
$VOT\,t^{ff}_{d,r}\cdot\texttt{operating\_cost\_scale}$，未服务变量系数为
$\max(0,VOT\,T)$。站点逐步功率和插头约束实际为
\begin{align}
 \sum_{d,r\ni s}p^{default}h_{d,r}-P_s^{default}z_s&\le0,\\
 \sum_{d,r\ni s}h_{d,r}-B_s&\le0.
\end{align}
依据：\srcpath{src/ev\_power\_traffic/infra\_design\_milp.cpp:35--201}。

若启用道路容量，G 并未精确积分 BPR；它创建 $f_a,e_a$，使用固定
$\alpha=0.15,\beta=4$ 和 $\alpha\beta=0.6$：
\begin{align}
 f_a&=\sum_{d,r:a\in r}h_{d,r},\\
 f_a-e_a&\le cap_a,\qquad f_a\le2cap_a,
\end{align}
并给 $f_a$、$e_a$ 分别赋 $VOT\,t_a^0\,scale$ 与其 $0.6$ 倍的线性系数。启用 LTM 动态行时，
这些静态路径/BPR 代理系数被清零，改由累计计数 VHT 系数承载。依据：
\srcpath{src/ev\_power\_traffic/infra\_design\_milp.cpp:203--288}、
\srcpath{src/ev\_power\_traffic/infra\_design\_milp.cpp:290--370}。

\subsection{H：滚动 LTM LP}
H 每个仿真步只执行当前滚动窗 LP 的第一个控制量。对站点，队列、能量和 V2G 行为为
\begin{align}
 Q_{j+1}-Q_j-b_j+c_j&=0,\qquad c_j\le Q_j,\\
 SOC_{j+1}-SOC_j-Eb_j+Ec_j+\Delta t\,v_j^{2g}&=0,\\
 v_j^{2g}-0.8SOC_j&\le0.
\end{align}
目标系数只对 $Q_{j+1}$ 加 \texttt{weight\_queue}，对各跟踪路段的
$N^{in}_{j+1}-N^{out}_{j+1}$ 加 \texttt{weight\_travel\_time}$\Delta t$，对 V2G 加
$-\texttt{weight\_v2g}$。LTM 行包括累计计数单调、发送/接收时滞、逐步容量，以及
$b_j\le N^{out}_{access,j+1}-N^{out}_{access,j}$。依据：
\srcpath{src/ev\_power\_traffic/ltm\_mpc.cpp:216--419}。

H 优先使用 HiGHS LP，失败后用本地单纯形；只有求解成功才推进持久队列、SOC 和累计计数。
\texttt{proven\_optimal} 还要求原始与整数残差均不超过 $10^{-7}$。依据：
\srcpath{src/ev\_power\_traffic/ltm\_mpc.cpp:423--514}。

\subsection{覆盖边界}
本章逐式覆盖 CTM 内部元胞、F 迭代、C 价格固定点、G 核心约束和 H 滚动行。D、E、E--LTM
规模较大且按模式分支，本章只定位其权威装配和证书语义；没有使用通用 CTM、Beckmann 或 MPEC
方程冒充未逐行转写的源码。
